\chapter{Review of the Literature}
\label{ch:litreview}

\draftnote{this chapter is a skeleton built only from the references already
collected in \texttt{references.bib} (all currently placeholder entries that
must be verified against the published record before submission). It is
deliberately scoped to the material needed to justify the unifying
hypothesis of \cref{ch:introduction} and the specific hypotheses tested in
\cref{ch:resolution,ch:tl-horizontal,ch:tl-vertical,ch:phase-plane,ch:fluidflow}
--- add literature on (1) time-lapse/4D geophysical monitoring, (2)
phase-based displacement sensing in other modalities (e.g.\ InSAR), and (3)
weighted least-squares plane-fitting precedents, none of which are yet
represented in the bibliography.}

\section{GPR Physics and Imaging}

Ground-penetrating radar images subsurface dielectric contrasts by recording
the reflections of a transmitted electromagnetic pulse; the underlying
electromagnetic principles, antenna behaviour, and propagation losses that
govern the achievable resolution and penetration depth of a survey are
covered in \citet{annan2009} and \citet{daniels2004}. Both texts establish
the half-wavelength rule of thumb for resolving two closely spaced
reflectors, which motivates the resolution-limit experiment of
\cref{ch:resolution}.

\section{Migration Algorithms}

Three families of migration algorithm are used throughout this thesis, and
each has an established lineage in the geophysical literature.
\citet{claerbout1985} introduces the exploding-reflector model that underlies
the zero-offset migration convention adopted in \cref{ch:theory}.
\citet{schneider1978} formulates Kirchhoff (integral, delay-and-sum)
migration as the adjoint of a forward modelling operator, the formulation
implemented here via the PyLops Kirchhoff operator. \citet{gazdag1978}
develops the frequency--wavenumber phase-shift migration algorithm used as
an independent wave-equation-based cross-check; \citet{stolt1978} is the
related constant-velocity Fourier-domain migration that the phase-shift
method generalises to depth-variable continuation. Together these three
algorithms make the resolution and time-lapse results of
\cref{ch:resolution,ch:tl-horizontal,ch:tl-vertical,ch:fluidflow}
independent of any single migration assumption.

\section{Simulation Tooling}

All synthetic data in this thesis are generated with gprMax
\citep{gprmax}, an open-source finite-difference time-domain electromagnetic
simulator purpose-built for GPR. Using a full-wave simulator rather than a
ray-based or convolutional forward model lets the back-propagation
(time-reversal) migration of \cref{ch:theory} be tested directly against the
same finite-difference physics used to generate the data, removing modelling
mismatch as a confound when comparing migration algorithms.

\section{Phase-Based Interpretation}

The use of signal \emph{phase}, rather than amplitude, as the primary
observable is the central methodological choice of this thesis.
\citet{castagna2016} demonstrates phase decomposition as a tool for
seismic interpretation, separating amplitude and phase information that is
otherwise conflated in a conventional seismic or GPR trace; the
phase-decomposition routines referenced in this thesis's processing pipeline
(\cref{ch:methodology}) build on this idea. This thesis extends phase-based
analysis from single-survey interpretation to a \emph{cross}-survey,
time-lapse setting, where the quantity of interest is not the phase of a
single image but the phase difference between two migrated images of the
same subsurface region.

\bigskip
\noindent The following chapter derives the migration and phase-plane theory
introduced qualitatively above in full, starting from the exploding-reflector
model and the three migration algorithms reviewed here, and building up to
the weighted least-squares phase-plane estimator validated in
\cref{ch:phase-plane}.
